
\section{Entrainment of bed sediment into still water}

A tank of still water 1~m deep stands over a fixed bed that is flat on one
half and slopes at $S = 10^{-3}$ on the other. The sediment is 0.5~mm sand
with the settling velocity $v_s$ of Ferguson and Church~\cite{FergusonChurch2004}.
There is no flow, so the case uses the depth--slope closure, under which the
bed shear stress of a cell depends only on its depth and bed slope,
$\tau_b = \rho g h S$, and the Shields stress is $\tau^{*} = hS/(RD)$ with
$R = \rho_s/\rho - 1$. On the flat half $\tau^{*} = 0$ lies below the critical
value $\tau_c^{*} = 0.04$ and nothing may happen: that half is the control.
On the sloped half $\tau^{*} \approx 1.2 \approx 30\,\tau_c^{*}$ and the
non-cohesive entrainment law of Smith and McLean fires,
\begin{equation}
E = v_s E^{*}, \qquad
E^{*} = \frac{0.65\,\gamma_0 X}{1 + \gamma_0 X}, \qquad
X = \frac{\tau^{*}}{\tau_c^{*}} - 1, \qquad \gamma_0 = 0.0024 ,
\end{equation}
while deposition returns sediment to the bed at the well-mixed rate
$D = d^{*} v_s c$ with $d^{*} = 1$. Every cell is a closed system, so the
sediment mass per unit area $m = hc$ obeys $dm/dt = E - D$, and from $c = 0$
the concentration relaxes to the equilibrium value $c_{eq} = E/(d^{*}v_s)$
along
\begin{equation}
c(t) = c_{eq}\left(1 - e^{-t/T}\right), \qquad T = \frac{h}{d^{*} v_s} \approx 10~\mathrm{s} .
\end{equation}
The bed is held fixed on purpose: with it evolving, the cells along the far
wall, whose reconstructed bed slope is zero, would not erode while their
neighbours did, and the steps that opened between them would feed back into
the depth--slope closure, so no cell would keep a known slope. The bed update
is checked by the settling case instead. The case checks the entrainment law
and its threshold, the depth--slope closure, and the balance of entrainment
against deposition in the conserved sediment mass. The solver reconstructs
the bed edge values each step; the limiter reduces the slope of the cells at
the kink and along the far wall, and those eight cells are compared against
the slope they were actually given.

\begin{figure}
\begin{center}
\includegraphics[width=0.9\textwidth]{concentration_plot.png}
\end{center}
\caption{Depth-averaged concentration against time, averaged over the sloped
half and over the flat (control) half.}
\end{figure}

\begin{figure}
\begin{center}
\includegraphics[width=0.9\textwidth]{profile_plot.png}
\end{center}
\caption{Concentration in every cell along the tank at the end of the run.}
\end{figure}

\endinput
